% ============================================================================
% Section N: Recorded Panel Discussion (Deliverable N)
% ============================================================================
\section{Panel Discussion}
\label{sec:N}


Table~\ref{tab:panel_plan} summarizes the 6-stage interactive discussion sequence aligned with the Practical~3 guidelines.

\begingroup
\footnotesize
\renewcommand{\arraystretch}{1.20}
\begin{xltabular}{\textwidth}{L{1.25cm} L{2.85cm} L{2.75cm} Y}
\caption{Structured 6--8-minute interactive panel discussion sequence for Group 27.}
\label{tab:panel_plan}\\
\toprule
\hdrrow \tblhead{Time} & \tblhead{Discussion Topic} & \tblhead{Lead \& Challengers} & \tblhead{Core Engineering Points \& Interactive Cross-Questions} \\
\midrule
\endfirsthead
\multicolumn{4}{l}{\footnotesize\color{SlateMuted}\textit{Table~\ref{tab:panel_plan} continued from previous page.}}\\
\toprule
\hdrrow \tblhead{Time} & \tblhead{Discussion Topic} & \tblhead{Lead \& Challengers} & \tblhead{Core Engineering Points \& Interactive Cross-Questions} \\
\midrule
\endhead
\bottomrule
\endfoot

0:00--1:10 & \textbf{1. System-of-Interest \& Boundary} & Lead: \memberOneName; Q\&A: Members 2 \& 3 & Why bound UA-EWRRF to a 10~t/d Nagpur catchment instead of a national network? Challenge: why exclude precious-metal pyrometallurgical smelting and Butibori TSDF from the boundary while retaining smart kiosks and ZLD? \\
\midrule
1:10--2:25 & \textbf{2. Concurrent Life Cycles} & Lead: \memberOneName; Q\&A: Member 3 & How LC-1 (plant), LC-2 (civil/skid build), LC-3 (spares/EHS), and LC-4 (decommissioning escrow) interlock from Phase~2. Challenge: what fails first if LC-3 blade stocking or LC-4 epoxy containment is deferred until Phase~6? \\
\midrule
2:25--3:40 & \textbf{3. DDPs vs.\ DIPs Classification} & Lead: \memberTwoName; Q\&A: Members 1 \& 3 & Hardest parameter to classify: \textit{Incoming WEEE Volume \& Mix (DIP-2)}---can kiosk incentives partially influence volume? Debate concludes total urban mix and OEM device gluing remain outside facility design authority, justifying DIP-2. \\
\midrule
3:40--4:55 & \textbf{4. Most Critical TPM \& Design Criterion} & Lead: \memberTwoName; Q\&A: Member 3 & Debate between \textbf{TPM-4 ($E_{\text{PM}}$)} (regulatory license-to-operate gate: $\le 30/2.0\text{ mg/m}^3$, MPCB standard) and \textbf{TPM-1 (MRR $\ge 82\%$)} (mission \& EPR revenue driver); explain benchmark vs.\ statutory basis. \\
\midrule
4:55--6:10 & \textbf{5. Design Trade-Off Selection} & Lead: \memberTwoName; Q\&A: Members 1 \& 3 & Why select \textbf{Alternative~2 (Modular Hybrid, 8.62/10)} over A1 and A3? Explicitly acknowledge accepted disadvantages: higher upfront CapEx than A1 (INR 7.8 vs.\ 3.2~Cr), lower peak burst speed than A3, and sensor calibration overhead. \\
\midrule
6:10--7:30 & \textbf{6. Changed DIP Condition \& Verdict} & Lead: \memberThreeName; Q\&A: Members 1 \& 2 & Show how Year-6 DIP-1/2/4/5 shifts (Li-ion/BFR mandate, +30\% power, +35\% TSDF fees) break unadapted A2 yet validate modular evolution; discuss when extreme labor scarcity + $>30\text{ t/d}$ volume would make automated pre-screened A3 preferable. \\

\end{xltabular}
\endgroup
